\documentclass[a4paper]{report} % Changé de 'article' à 'report'

%% Language and font encodings
\usepackage[english, french]{babel}
\usepackage[utf8]{inputenc}
\usepackage[T1]{fontenc}

%% Sets page size and margins
\usepackage{geometry}
\usepackage{lipsum}              % Pour du texte d'exemple (facultatif)

%% Useful packages
\usepackage{amsmath, amssymb}
\usepackage{amssymb} % Redundant if amsmath is used, but kept as in original
\usepackage{graphicx}
\usepackage[colorlinks=true, allcolors=black]{hyperref}
\usepackage{float}
\usepackage{tikz}
\usepackage{calc}
\usepackage{pgfplots}
\usepackage{setspace}
\usepackage{titlesec}
\usepackage{parskip}
\usepackage{tabularx}
\usepackage{booktabs}
\usepackage{enumitem}
\usepackage{xcolor}
\definecolor{accenture}{HTML}{A100FF}
\usetikzlibrary{positioning, shapes.geometric, arrows.meta, backgrounds}
\usetikzlibrary{shapes, arrows, positioning}
\setlist[itemize]{itemsep=8pt, topsep=2pt} % Applique ces réglages à toutes les listes itemize
\setcounter{secnumdepth}{3} %pour que le niveau 3 (subsubsection soit numeroté)
\setcounter{tocdepth}{2} %pour que le niveau 2 (subsection apparaisse dans la table des matieres)
% Personnalisation des titres de chapitres
% Met le mot "Chapitre" et son numéro en gras et très grand (\Huge)
% et le titre du chapitre également en \Huge
% Redéfinition complète du format des pages de chapitre


% --- 1. Apparence des titres (police, taille, gras...) ---

\titleformat{\chapter}[display]
  {\normalfont\Huge\bfseries} % Apparence générale : gras et centré
  {\vspace*{50pt}\chaptertitlename\ \thechapter} % Espace au-dessus + Texte "Chapitre X"
  {25pt} % Espace vertical entre "Chapitre X" et le titre du chapitre
  {\Huge} % Taille de la police pour le titre du chapitre
  [\vfill\newpage] % Reste de la page vide et saut de page

\titleformat{\section}
  {\cleardoublepage\normalfont\LARGE\bfseries} % Police grande et en gras
  {\thesection}
  {1em}
  {}

\titleformat{\subsection}
  {\normalfont\large\bfseries} % Police "large" et en gras
  {\thesubsection}
  {1em}
  {}

\titleformat{\subsubsection}
  {\normalfont\normalsize\bfseries} % Police normale et en gras
  {\thesubsubsection}
  {1em}
  {}


% --- 2. Espacement avant et après les titres ---

% Syntaxe : \titlespacing*{commande}{gauche}{AVANT}{APRES}

\titlespacing*{\section}
  {0pt}      % gauche
  {10ex}      % avant (l'espace que vous vouliez ajouter)
  {3ex}    % après

\titlespacing*{\subsection}
  {0pt}      % gauche
  {6ex}    % avant
  {2.5ex}    % après

\titlespacing*{\subsubsection}
  {0pt}      % gauche
  {4ex}   % avant
  {2.5ex}    % après

\pgfplotsset{compat=1.14}

\setlength{\intextsep}{16pt} % Espace autour des figures dans le texte

%% Added packages for new layout
\usepackage{fancyhdr} % For custom headers and footers

\newcommand\fillin[1]{\makebox[#1]{\dotfill}}


\begin{document}
% La page de garde
\begin{titlepage} % Utiliser l'environnement titlepage pour la page de garde
% Définir une géométrie spécifique pour la page de titre pour éviter le débordement
\newgeometry{inner=1cm, outer=1cm, top=1cm, bottom=0.5cm, includefoot=false, includehead=false}
\pagenumbering{roman} % La numérotation commence ici si désiré pour les pages liminaires
\thispagestyle{empty} % La page de garde ne doit pas avoir d'en-tête/pied de page

% --- START OF FIRST PAGE CONTENT (Logos d'origine rétablis) ---
\begin{tikzpicture}[remember picture, overlay]
% IMPORTANT: Ensure these image files exist in the 'images' subfolder and are valid!
% Logos d'origine de la première page
\node[xshift=-7cm, yshift=-3.1cm] at (current page.north) {\includegraphics[scale=0.11]{images/logo_ida.png}};
\node[below = 0.5cm] at (current page.north) {\includegraphics[scale = 0.11]{images/ucbl.jpg}};
\node[xshift=6.8cm, yshift=-3.3cm] at (current page.north) {\includegraphics[scale=0.11]{images/isfa_logo.jpg}};
\end{tikzpicture}

\vspace*{2.5cm} % Réduction de l'espace vertical

{\Large
{\bfseries
\begin{center}
Mémoire présenté le : \\[0.3cm]
pour l'obtention du Diplôme Universitaire d'actuariat de l'ISFA \\
et l'admission à l'Institut des Actuaires
\end{center}
}
\vspace*{0.5cm} % Réduction de l'espace vertical

\noindent Par : Gaylord LEGRIS \\[0.3cm]
\noindent Titre : Analyse de Sensibilités sur des Portefeuilles de Passifs en Assurance Vie selon les indicateurs Solvabilité 2 : Approche par Générations et Agrégations des portefeuilles
\\[0.3cm]
\noindent Confidentialité : \quad $\boxtimes$ NON \qquad $\square$ (Durée : $\square$ 1 an \quad  $\square$ 2 ans)
\\
\textit{Les signataires s'engagent à respecter la confidentialité indiquée ci-dessus}
\vspace*{0.5cm}

\noindent
\begin{minipage}{0.6\textwidth}
\begin{tabular}{p{7cm}p{4cm}}
\textit{Membres présents du jury de l'Institut des Actuaires} & \textit{Signature} \\[1.2cm]
\fillin{6cm} & \\[0.8cm]
\fillin{6cm} & \\[0.8cm]
\fillin{6cm} & \\[0.8cm]
\textit{Membres présents du jury de l'ISFA} & \\[1.2cm]
\fillin{6cm} & \\[0.8cm]
\fillin{6cm} & \\[0.8cm]
\fillin{6cm} & \\[0.8cm]
\end{tabular}
\rule{0mm}{4.6cm} % Ensure minipage height
\end{minipage}%
\begin{minipage}{0.4\textwidth}
\textit{Entreprise :} \\
\textit{Nom :} \\[0.3cm]
\textit{Signature :} \\[0.5cm]
\textit{Directeur de mémoire en entreprise :} \\
\textit{Nom :}  \\[0.3cm]
\textit{Signature :} \\[0.5cm]
\textit{Invité :} \\
\textit{Nom :} \\[0.3cm]
\textit{Signature :} \\[0.5cm]
\textit{\textbf{Autorisation de publication et de mise en ligne sur un site de diffusion de documents actuariels} (après expiration de l'éventuel délai de confidentialité)}\\[0.2cm]
Signature du responsable entreprise \\
\framebox[7cm]{\rule{0mm}{1.5cm}}
Signature du candidat \\
\framebox[7cm]{\rule{0mm}{1.5cm}}
\end{minipage}
}
% --- END OF FIRST PAGE CONTENT ---
\end{titlepage}
\restoregeometry % Restaurer la géométrie précédente avant de la redéfinir pour le reste du document.

% Redéfinir la géométrie pour le reste du document après la page de garde
\newgeometry{left=2.54cm,right=2.54cm,top=3.54cm,bottom=3.54cm}
\onehalfspacing
% Augmenter headheight pour accommoder des logos plus grands
\setlength{\headheight}{50pt} % Ajusté pour des logos de 1.6cm de hauteur

% Configurer fancyhdr pour les en-têtes et pieds de page personnalisés
\pagestyle{fancy}
\fancyhf{} % Effacer tous les champs d'en-tête et de pied de page

% En-têtes des pages suivantes avec les logos demandés (taille doublée)
% IMPORTANT: Assurez-vous que images/accenture_logo.png existe !
\fancyhead[L]{\includegraphics[height=1.6cm]{images/accenture_logo.png}} % Logo Accenture à gauche
\fancyhead[R]{\includegraphics[height=1.6cm]{images/isfa_logo.jpg}}   % Logo ISFA à droite

\fancyfoot[C]{\sffamily \thepage}

\renewcommand{\headrulewidth}{0pt}
\renewcommand{\footrulewidth}{0pt}

% Redéfinir le style 'plain' pour qu'il soit identique à 'fancy' afin que les logos apparaissent sur les premières pages de chapitre (taille doublée)
\fancypagestyle{plain}{%
  \fancyhf{}%
  \fancyhead[L]{\includegraphics[height=1.6cm]{images/accenture_logo.png}}% Logo Accenture à gauche
  \fancyhead[R]{\includegraphics[height=1.6cm]{images/isfa_logo.jpg}}%   Logo ISFA à droite
  \fancyfoot[C]{\sffamily \thepage}%
  \renewcommand{\headrulewidth}{0pt}%
  \renewcommand{\footrulewidth}{0pt}%
}


% Contenu principal en une seule colonne
\tableofcontents

% IMPORTANT: Ensure all included .tex files exist in the specified paths
% and do not contain LaTeX errors themselves!
\chapter*{Résumé}
% Ajouter manuellement cette section à la table des matières
\addcontentsline{toc}{chapter}{Résumé}
Petit résumé en français de mon mémoire ?
\chapter*{Abstract}
% Ajouter manuellement cette section à la table des matières
\addcontentsline{toc}{chapter}{Abstract}

It's a brief sum up in english !
\chapter*{Remerciements}
% Ajouter manuellement cette section à la table des matières
\addcontentsline{toc}{chapter}{Remerciements}

Pour la réalisation de ce mémoire, je tiens à adresser mes remerciements à toutes les personnes qui ont contribué de près ou de loin à son élaboration. Que ce soit au travail avec Lionel Aldeberd qui m'a appris énormément, Lucas Blancheton qui a toujours été là pour m'aiguiller dans mes choix. Merci à Luc pour les relectures et à tous mes collègues stagiaires et alternants, Antoine, Cindy, Eleonore, Lucile, Malak, Manon, Nicolas, Tanguy et Titouan pour leur bonne humeur ! J'ai vraiment passé des très chouettes moments à vos côtés !
\chapter*{Synthèse}
% Ajouter manuellement cette section à la table des matières
\addcontentsline{toc}{chapter}{Synthèse}

Un long résumé en français de mon mémoire 
\chapter*{Synthesis}
% Ajouter manuellement cette section à la table des matières
\addcontentsline{toc}{chapter}{Synthesis}

Un long résumé en anglais de mon mémoire
\chapter*{Introduction}
% Ajouter manuellement cette section à la table des matières
\addcontentsline{toc}{chapter}{Introduction}

Placement privilégié des épargnants français, l'assurance vie a atteint un encours record de 1 923 milliards d'euros à fin 2023 (France Assureurs), confirmant ainsi son rôle prépondérant dans le patrimoine financier national. Toutefois, ce secteur fait face à une rupture structurelle marquée par la fin du cycle de taux bas et la remontée brutale des taux d'intérêt observée depuis 2022. Le taux de revalorisation moyen des fonds en euros a ainsi atteint 2,6 \% pour l'année 2023 (estimation ACPR), créant un paradigme nouveau. Cette mutation rend d'autres produits d'épargne plus attractifs et exerce une pression concurrentielle inédite sur les contrats d'assurance vie, notamment sur les fonds en euros qui subissent une décollecte nette significative (\textbf{pas sûr, faudra que je regarde la replay de l'analyse SFCR fait ya pas longtemps}). Pour les assureurs, le défi est de taille : leurs portefeuilles d'actifs, majoritairement constitués d'obligations acquises durant la longue période de taux bas, présentent une forte inertie. Cet héritage obligataire freine leur capacité à servir des rendements compétitifs et place la gestion actif-passif (ALM) au cœur des enjeux stratégiques.\newline{}

Pour piloter leur bilan, les assureurs s'appuient sur des modèles ALM sophistiqués, essentiels pour simuler l'impact de différentes stratégies dans le cadre réglementaire de Solvabilité II. Cependant, la complexité de ces modèles et la nécessité de réaliser un grand nombre de simulations se heurtent à une contrainte opérationnelle majeure : le temps de calcul. Cette contrainte limite la capacité des assureurs à explorer en profondeur l'ensemble des risques et des opportunités. Face à cette réalité, une question centrale émerge : dans quelle mesure une agrégation des engagements de passif permet-elle de préserver la fidélité des indicateurs de risque tout en optimisant les temps de calcul ? Ce mémoire se propose d'investiguer cette problématique en étudiant l'impact de l'agrégation des portefeuilles de passifs. L'enjeu est de déterminer si une représentation plus grossière du passif peut suffire pour le pilotage stratégique et sous quelles conditions une telle simplification est valide, sans masquer des dynamiques de risque essentielles.\newline{}

Pour répondre à cette problématique, ce mémoire adoptera une double approche. Premièrement, nous reconnaissons l'hétérogénéité des engagements, qui ont été souscrits dans des contextes de taux variés. Une approche par générations de passifs est donc essentielle pour distinguer les dynamiques propres à chaque cohorte de contrats. Cette démarche de génération de portefeuilles synthétiques est fondamentale pour tester en amont l'attractivité et la résilience de nouvelles offres.\newline{}

Deuxièmement, l'analyse portera sur les effets de l'agrégation de ces différentes générations de passifs. L'objectif est de comprendre comment les risques se combinent et si des effets de portefeuille permettent une mutualisation. Ce mémoire ne se contentera pas d'analyser l'impact d'une seule méthode d'agrégation ; au contraire, nous testerons et comparerons plusieurs approches. Le critère de sélection de la méthode la plus pertinente reposera sur un triple objectif : minimiser l'écart des indicateurs clés (notamment le Best Estimate et le SCR), optimiser la rapidité des calculs et atteindre le plus haut niveau d'agrégation possible qui garderait une significaitivité économique pour l'assureur suffisante.\newline{}

L'axe principal de ce mémoire consistera donc à mener une analyse de sensibilités approfondie sur ces portefeuilles, qu'ils soient granulaires ou agrégés. Notre étude s'appuiera sur des indicateurs quantitatifs clés issus de la norme Solvabilité II, en évaluant notamment l'impact des chocs économiques sur le Solvency Capital Requirement (SCR) et la Present Value of Future Profits (PVFP) (\textbf{Je ne sais pas encore si je regarderai que ces indicateurs, je pense aussi regarder l'impact de chocs sur certains colonnes de l'input sur ce genre d'indicateurs aussi}). Ces métriques permettront de mesurer rigoureusement comment l'agrégation modifie la perception du risque et la valeur économique du portefeuille.\newpage{}
Ce mémoire s'articulera en quatre parties distinctes, chacune conçue pour apporter une réponse progressive et rigoureuse à notre problématique.\newline{}


La \textbf{première partie} sera consacrée au cadre conceptuel de notre étude. Nous y détaillerons le contexte réglementaire de Solvabilité II, qui définit les exigences de capital et les métriques de risque, ainsi que les principes fondamentaux de la modélisation actif-passif (ALM). L'avantage de cette section est de fournir au lecteur les clés de compréhension essentielles pour appréhender les enjeux techniques et stratégiques du pilotage d'un bilan assurantiel.\newline{}


La \textbf{deuxième partie} adoptera une approche pratique en se concentrant sur la mise en œuvre de notre environnement de simulation. Nous y décrirons la méthodologie de création des portefeuilles de passifs synthétiques, représentatifs de différentes générations de contrats, ainsi que les outils développés pour leur projection. Cette étape est cruciale car elle garantit la robustesse et la pertinence des analyses qui suivront, en créant un laboratoire d'expérimentation fiable.\newline{}


La \textbf{troisième partie} constituera le cœur méthodologique de ce mémoire. Elle explorera et comparera de manière systématique plusieurs techniques d'agrégation des engagements de passif. L'objectif sera d'identifier les approches les plus prometteuses, en évaluant leur capacité à simplifier la structure du portefeuille sans dénaturer ses caractéristiques fondamentales. Cette analyse comparative permettra de mettre en lumière les forces et faiblesses de chaque méthode.\newline{}


Enfin, la \textbf{quatrième partie} présentera et analysera les résultats de nos simulations. À travers des tests de sensibilité approfondis sur les indicateurs clés (SCR, PVFP), nous quantifierons l'impact de chaque méthode d'agrégation sur la perception du risque et la valeur économique. Cette analyse empirique nous permettra de conclure sur la validité des approches testées et de formuler des recommandations concrètes sur les conditions d'utilisation d'un passif agrégé pour un pilotage ALM à la fois efficace et optimisé.

\clearpage
\pagenumbering{arabic}

\chapter{Introduction au contexte réglementaire et à la modélisation ALM}


% Fichier: 1_1_specificites_assurance_vie.tex
% Section: 1.1 Spécificités de l'assurance vie

\section{Spécificités de l'assurance vie}
\label{sec:specificites_assurance_vie}

L'assurance vie demeure un pilier central de la stratégie patrimoniale des Français. En effet, elle représente le placement le plus important en valeur, totalisant 1989 milliards d’euros d’encours sur les différents contrats à fin décembre 2024. Avec un capital moyen excédant 100\,000 euros par souscripteur, son rôle prépondérant dans l'épargne nationale est indéniable. Cette popularité s'explique par une combinaison unique d'utilités et de caractéristiques qui façonnent sa gestion et sa modélisation. Comprendre ces spécificités est fondamental avant d'aborder le cadre réglementaire de Solvabilité II et les subtilités de la modélisation Actif-Passif (ALM) qui sont l'objet de ce mémoire.

\subsection{Utilité de l'assurance vie}
\label{subsec:utilite_assurance_vie}
L'assurance vie répond à une multitude d'objectifs patrimoniaux, ce qui explique son attrait constant :
\begin{itemize}
    \item \textbf{Constitution d'une épargne à long terme :} C'est un outil privilégié pour se constituer un capital progressivement, en vue de financer des projets futurs (retraite, études des enfants, acquisition immobilière) ou simplement pour valoriser un capital existant.
    \item \textbf{Transmission du patrimoine :} L'assurance vie offre un cadre fiscal avantageux pour la transmission d'un capital à des bénéficiaires désignés, que ce soit en ligne directe ou non. Les abattements spécifiques et la possibilité de désigner librement les bénéficiaires en font un instrument de planification successorale efficace.
    \item \textbf{Protection des proches :} En cas de décès de l'assuré, le versement d'un capital ou d'une rente aux bénéficiaires permet de les protéger financièrement et de faire face aux conséquences pécuniaires de la disparition.
    \item \textbf{Recherche de rendement et diversification :} Avec une gamme de supports allant des \textbf{fonds en euros} sécurisés aux \textbf{unités de compte (UC)} plus dynamiques, l'assurance vie permet d'adapter son investissement à son profil de risque et à ses objectifs de rendement, tout en offrant des possibilités de diversification.
    \item \textbf{Avantages fiscaux :} Outre la transmission, l'assurance vie bénéficie d'une fiscalité allégée sur les plus-values en cas de rachat après une certaine durée de détention (notamment après 8 ans), ce qui renforce son attractivité comme placement de long terme.
\end{itemize}
Cette polyvalence fait de l'assurance vie un produit incontournable, mais engendre également une complexité dans sa gestion pour les assureurs.

\subsection{Dualité du produit : Épargne et Prévoyance}
\label{subsec:dualite_produit}
Au-delà de ses utilités multiples, l'une des premières spécificités de l'assurance vie réside dans sa nature hybride. Elle combine à la fois une dimension d'épargne à long terme et une composante de prévoyance. 
\begin{itemize}
    \item \textbf{Volet Épargne :} Les contrats d'assurance vie permettent aux souscripteurs de se constituer un capital sur la durée, en bénéficiant potentiellement de rendements attractifs et d'un cadre fiscal avantageux, notamment en cas de rachat après une certaine période de détention ou pour la transmission du capital.
    \item \textbf{Volet Prévoyance :} En cas de décès de l'assuré, le capital constitué (ou un capital garanti) est versé à un ou plusieurs bénéficiaires désignés, leur offrant ainsi une protection financière. Cette dimension de couverture du risque de décès est inhérente à de nombreux contrats.
\end{itemize}
Cette dualité engendre des flux financiers complexes et des engagements de long terme pour l'assureur, nécessitant une gestion prudente et prospective. C'est pourquoi il met en place des moyens de calculer les montants qu'il doit

\subsection{La garantie du capital et le taux minimum garanti (TMG)}
\label{subsec:garantie_capital_tmg}
Historiquement, de nombreux contrats d'assurance vie en euros se sont distingués par la \textbf{garantie du capital} investi par l'assuré. L'assureur s'engage à restituer au minimum les sommes versées, nettes de frais. 
À cela s'ajoute souvent la notion de \textbf{Taux Minimum Garanti (TMG)}, qui est un rendement plancher que l'assureur s'engage à servir annuellement sur le capital. Bien que la tendance réglementaire et les conditions de marché aient conduit à une baisse, voire une disparition des TMG élevés sur les nouveaux contrats, cette caractéristique a un impact majeur sur les engagements passés et la gestion des actifs en couverture. Les assureurs doivent en effet générer des rendements financiers suffisants pour honorer ces garanties, ce qui influence directement leur politique d'investissement et leur exposition aux risques de marché.

\subsection{La participation aux bénéfices (PB)}
\label{subsec:participation_benefices}
Au-delà du TMG, les assurés bénéficient de la \textbf{participation aux bénéfices (PB)}. Les assureurs ont l'obligation de redistribuer une partie des bénéfices techniques et financiers qu'ils réalisent sur la gestion des contrats d'assurance vie. Cette redistribution vient s'ajouter au TMG pour former le rendement global servi à l'assuré. La politique de PB est un levier important pour l'assureur, lui permettant de lisser les rendements dans le temps et de gérer les attentes des assurés, tout en respectant les contraintes réglementaires (provision pour participation aux excédents - PPE). La gestion de la PB est un élément central de la stratégie ALM.

\subsection{Liquidité et faculté de rachat}
\label{subsec:liquidite_rachat}
Les contrats d'assurance vie offrent généralement une grande \textbf{liquidité} aux souscripteurs, qui disposent d'une \textbf{faculté de rachat} total ou partiel de leur épargne à tout moment (souvent après une période initiale et sous conditions fiscales). Cette caractéristique, bien qu'avantageuse pour l'assuré, représente un risque pour l'assureur :
\begin{itemize}
    \item \textbf{Risque de rachat massif (risque de ``bank run'' ou de ``surrender'') :} En cas de crise de confiance, de forte hausse des taux d'intérêt alternatifs ou d'autres événements défavorables, les assureurs pourraient faire face à une vague de rachats importante, les contraignant à liquider des actifs dans des conditions potentiellement défavorables.
    \item \textbf{Antisélection :} Les comportements de rachat peuvent dépendre des conditions de marché et de la situation individuelle des assurés, introduisant un biais dans les flux de sortie.
\end{itemize}
La modélisation des comportements de rachat est donc un enjeu crucial pour l'évaluation des passifs et la gestion ALM.

\subsection{Horizon de gestion à long terme}
\label{subsec:horizon_long_terme}
Les engagements en assurance vie s'inscrivent typiquement sur un \textbf{horizon de long, voire très long terme}. Les contrats peuvent durer plusieurs décennies, et les prestations (rachats, décès) s'étalent dans le temps. Cette perspective temporelle impose aux assureurs :
\begin{itemize}
    \item Une adéquation entre la duration des actifs et celle des passifs.
    \item Une prise en compte des évolutions futures des taux d'intérêt, de l'inflation, de la mortalité et des comportements des assurés.
    \item Une capacité à anticiper les changements réglementaires et économiques sur de longues périodes.
\end{itemize}
Cette caractéristique est au cœur de la problématique ALM et justifie l'utilisation de modèles de projection stochastiques.

\subsection{Diversité des supports : Fonds en euros et Unités de Compte (UC)}
\label{subsec:fonds_euros_uc}
Le marché de l'assurance vie se segmente principalement entre :
\begin{itemize}
    \item \textbf{Les fonds en euros :} Majoritairement investis en obligations, ils offrent une garantie du capital et un rendement (TMG + PB). Le risque financier est principalement porté par l'assureur.
    \item \textbf{Les unités de compte (UC) :} Investies sur des supports plus dynamiques (actions, immobilier, OPCVM diversifiés), la valeur des UC fluctue en fonction des marchés financiers. Le risque financier est ici porté par l'assuré, l'assureur garantissant un nombre de parts et non une valeur en euros.
\end{itemize}
La part croissante des UC dans les nouveaux contrats modifie le profil de risque des assureurs et des assurés, et complexifie la modélisation des passifs, notamment en ce qui concerne les garanties optionnelles (garanties plancher en cas de décès sur UC, par exemple) et les frais prélevés. Ces deux types de support impliquent des passifs avec des caractéristiques et des risques très différents pour l'assureur.

\bigskip % Ajoute un peu d'espace vertical avant la transition vers la section suivante

Ces spécificités, interdépendantes, créent un environnement complexe pour les assureurs vie. Elles soulignent l'importance d'une modélisation ALM robuste et d'une analyse fine des sensibilités aux différents facteurs de risque, notamment dans le cadre exigeant de Solvabilité II que nous aborderons dans la section suivante.


\section{La réglementation Solvabilité 2}

La réglementation Solvabilité 2 (S2) est un cadre prudentiel harmonisé au niveau de l'Union Européenne, dont l'objectif principal est de renforcer la protection des assurés et bénéficiaires, d'accroître la compétitivité des assureurs européens et de mettre en place une évaluation économique du bilan des entreprises d'assurance et de réassurance. Entrée en vigueur le 1er janvier 2016, elle vise à assurer la pérennité et la solvabilité des organismes face aux risques inhérents à leur activité.

\subsection{Différents piliers}
La directive Solvabilité 2 repose sur trois piliers interdépendants couvrant des sujets liés à la quantification des risques, à la gouvernance et à la transparence.

\subsubsection{Pilier 1 : Exigences quantitatives}
Le Pilier 1 définit les exigences de fonds propres que les assureurs et réassureurs doivent détenir pour couvrir les risques liés à leurs opérations. Il détaille la méthodologie de calcul du capital réglementaire, ainsi que l'évaluation des actifs et des passifs en valeur de marché. Les provisions techniques, par exemple, sont définies comme la combinaison du Best Estimate (BE) et de la Risk Margin (RM). Le calcul du SCR (Solvency Capital Requirement) est également au cœur de ce pilier.
\textbf{IL FAUT QUE JE PARLE DAVANTAGE DE CETTE PARTIE CAR ELLE EST SUPER IMPORTANTE POUR LA SUITE DU MEMOIRE ET POUR SON UTILITE EN ALM}
\subsubsection{Pilier 2 : Exigences qualitatives}
Le Pilier 2 complète les exigences quantitatives par des exigences qualitatives en matière de gestion des risques et de gouvernance. Il requiert notamment la mise en place d'un système de gouvernance clair et bien défini, des mécanismes de contrôle interne adaptés à l'activité, et des processus efficaces pour évaluer, gérer et reporter les risques. L'évaluation interne des risques et de la solvabilité (ORSA - Own Risk and Solvency Assessment) constitue la pièce maîtresse de ce dispositif, visant à une identification globale et coordonnée des risques auxquels l'entreprise est exposée.

\subsubsection{Pilier 3 : Transparence et reporting}
Le Pilier 3 concerne la communication externe de l'assureur, imposant la transparence au travers de la publication de rapports qualitatifs (SFCR - Solvency and Financial Conditions Report, et RSR - Regulatory Solvency Report) et de rapports quantitatifs (QRT - Quantitative Reporting Templates) destinés au superviseur et parfois au public.

\subsection{Formule interne vs formule standard}
Pour le calcul du Solvency Capital Requirement (SCR), la réglementation Solvabilité 2 offre aux assureurs deux approches principales : la formule standard et le modèle interne.

\subsubsection{Formule standard}
La formule standard est une méthode prédéfinie par les régulateurs (EIOPA) qui prend en compte divers sous-modules de risques tels que le risque de marché, le risque de souscription en assurance vie et non-vie, le risque de contrepartie, etc.. Le SCR est calculé en appliquant des chocs de stress sur le scénario central pour mesurer l'impact sur le bilan (NAV). Les différents sous-modules de SCR sont agrégés à l'aide d'une matrice de corrélation fournie par le régulateur. Cette méthode est privilégiée car elle est la plus représentative de la globalité du marché.

\subsubsection{Formule interne}
Avec l'approbation des régulateurs, les entreprises peuvent développer leurs propres modèles statistiques et économétriques pour évaluer leur SCR, en fonction de leur profil de risque spécifique. Cette approche permet une mesure plus fine des risques et une allocation de capital potentiellement plus optimisée, car elle est adaptée aux spécificités de l'entreprise.

\subsubsection{Comparaison des deux approches}
Le choix entre la formule standard et le modèle interne dépend de la complexité de l'entreprise, de son profil de risque et de sa capacité à développer et à maintenir un modèle interne robuste et validé par les autorités de contrôle. La formule standard offre une approche simplifiée et harmonisée, tandis que le modèle interne permet une personnalisation et une optimisation plus poussées de l'exigence en capital. Toutefois, le développement et la validation d'un modèle interne sont des processus lourds et coûteux.

\section{Définitions et enjeux de l’ALM}
\subsection{Définition de l’ALM}
\subsection{Enjeux pour les assureurs}

\section{Présentation du modèle ALM (pour plus tard)}

\section{Les générateurs de scénarios économiques}
\subsection{Définition et rôle}
\subsection{Exemples de générateurs utilisés}

\section{Qu’est ce qu’un model point et pourquoi on les utilise ?}
\subsection{Définition des model points}
\subsection{Utilisation dans la modélisation ALM}

\section{Impact des Réglementations sur les Portefeuilles (A mettre plus tard)}
\subsection{Analyse de l'impact des réglementations sur les structures de portefeuilles de passifs}
\subsection{Quelles sont les réglementations existantes concernant l’agrégation en MP}
\subsection{Études de cas illustrant les contraintes réglementaires}



\chapter{Introduction au contexte réglementaire et à la modélisation ALM}
Ceci est le contenu du premier chapitre. Vous pouvez écrire ici votre texte ou ajouter des sections et sous-sections.

\section{Spécificités de l'assurance vie}

\section{La réglementation Solvabilité 2}
\subsection{Différents piliers}
\subsubsection{Pilier 1 : Exigences quantitatives}
\subsubsection{Pilier 2 : Exigences qualitatives}
\subsubsection{Pilier 3 : Transparence et reporting}

\subsection{Formule interne vs formule standard}
\subsubsection{Formule standard}
\subsubsection{Formule interne}
\subsubsection{Comparaison des deux approches}

\section{Définitions et enjeux de l’ALM}
\subsection{Définition de l’ALM}
\subsection{Enjeux pour les assureurs}

\section{Présentation du modèle ALM (pour plus tard)}

\section{Les générateurs de scénarios économiques}
\subsection{Définition et rôle}
\subsection{Exemples de générateurs utilisés}

\section{Qu’est ce qu’un model point et pourquoi on les utilise ?}
\subsection{Définition des model points}
\subsection{Utilisation dans la modélisation ALM}

\section{Impact des Réglementations sur les Portefeuilles (A mettre plus tard)}
\subsection{Analyse de l'impact des réglementations sur les structures de portefeuilles de passifs}
\subsection{Quelles sont les réglementations existantes concernant l’agrégation en MP}
\subsection{Études de cas illustrant les contraintes réglementaires}



\chapter{Contraintes techniques et création d'outils}
\section{Nécessité du générateur de portefeuille passif}
\subsection{Besoin de générer des données pour simuler un nouveau produit}
\subsection{Simuler différentes évolutions du business mix pour orienter politique de souscription/politique commerciale.}
\subsection{Simuler un portefeuille représentatif du marché ou composé des principaux concurrents pour se positionner.}
\section{Contraintes techniques associées}
\subsection{Mise à jour sur des outils plus récents}
\subsection{langage open source, permet de s'écarter des problématiques financières (coût de licence).}
\subsection{Travail sur des outils mis à jour fréquemment (Python, Polars).}

\section{Développement du Modèle ALM en Python}
\subsection{Présentation du modèle ALM développé pour Accenture}
\subsection{Fonctionnement du modèle et apprentissage personnel.}
\subsection{Limites du modèle à l’heure actuelle.}

\section{Générateur de portefeuille de passif}
\subsection{Description des contraintes techniques rencontrées}
Description des contraintes techniques rencontrées dans la génération des portefeuilles.

\subsection{Description du générateur de modèle point}
\subsection{Importance pour un cabinet de conseil et rassurance des clients.}
\chapter{Protocole d'Analyse : Sélection d'une Méthode d'Agrégation}

\section{Présentation des Méthodes d'Agrégation candidates}
    \subsection{Approches par clustering (K-means, DBSCAN/HDBSCAN)}
    % Votre texte ici...
    \subsection{Autres approches (MP par âge, MP Amine, etc.)}
    % Votre texte ici...

\section{Définition du Protocole de Test Comparatif}
    \subsection{Constitution des portefeuilles de test}
    % Votre texte ici...
    \subsection{Définition des critères de sélection : fidélité des indicateurs (BE/SCR), performance et temps de calcul}
    % Votre texte ici...

\section{Analyse Comparative et Choix de la Méthode Optimale}
    \subsection{Synthèse des performances pour chaque méthode candidate}
    % Votre texte ici...
    \subsection{Justification du choix de la méthode retenue pour l'analyse de sensibilité}
    % Votre texte ici...
% \chapter{Protocole d'Analyse : Agrégation et Scénarios de Sensibilité}

% % Introduction du chapitre : Expliquer que ce chapitre pose toute la méthodologie
% % qui sera appliquée dans la partie suivante. C'est le "comment" de l'analyse.

% \section{Méthodologie et Sélection du Modèle d'Agrégation}
% % Objectif : Décrire le processus complet qui mène au choix d'UNE méthode d'agrégation.

% \subsection{Description des Méthodes d'Agrégation Candidates}
% % Reprendre votre description technique des méthodes (K-means, DBSCAN, etc.)

% \subsection{Tests de Performance et Analyse comparative}
% % Fusionner la présentation des portefeuilles de test et l'analyse
% \subsubsection{Présentation des portefeuilles de test}
% % Décrire les portefeuilles utilisés spécifiquement pour comparer les méthodes d'agrégation.

% \subsubsection{Critères d'évaluation et résultats des tests}
% % Analyser les résultats (BE, SCR, temps de calcul) pour chaque méthode.
% % Inclure ici la réflexion sur l'optimisation du nombre de Model Points.

% \subsection{Choix du Modèle d'Agrégation Optimal pour les métriques S2}
% % Conclure cette section en justifiant le choix d'UNE méthode spécifique
% % au regard des résultats précédents et de sa compatibilité avec les architectures modernes.

% \section{Définition des Scénarios de Sensibilité}
% % Objectif : Décrire précisément les tests qui seront menés sur les portefeuilles.

% \subsection{Création des Portefeuilles de Test via le Générateur}
% % Expliquer comment le générateur est utilisé pour créer les portefeuilles de base
% % ainsi que leurs variations pour les tests de sensibilité.

% \subsection{Description des Chocs et Modifications Appliqués}
% % Détailler les chocs (positifs/négatifs) sur les variables clés.
% % Décrire le scénario d'ajout d'un nouveau produit (caractéristiques, volume, etc.).

\chapter{Analyse de Sensibilité des Indicateurs S2 aux Caractéristiques du Passif}

\section{Définition des Scénarios de Sensibilité}
    \subsection{Création des portefeuilles de test via le générateur}
    % Votre texte ici...
    \subsection{Description des chocs sur les variables clés (âge, montant de la PM, etc.)}
    % Votre texte ici...
    \subsection{Scénario d'intégration d'un nouveau produit dans le portefeuille}
    % Votre texte ici...

\section{Analyse de l'Impact de l'Agrégation sur la Mesure des Chocs}
    \subsection{Comparaison des indicateurs S2 sur portefeuilles choqués granulaires et agrégés}
    % Votre texte ici...
    \subsection{Analyse de la fidélité de la méthode d'agrégation à retranscrire la sensibilité}
    % Votre texte ici...

\section{Interprétation des Résultats et Validation de l'Approche}
    \subsection{Validation de la performance de la chaîne de modélisation (Générateur - Agrégation - Modèle ALM)}
    % Votre texte ici...
    \subsection{Enseignements sur la sensibilité des portefeuilles aux modifications de caractéristiques du passif}
    % Votre texte ici...

% \chapter{Résultats et Analyse des Sensibilités}

% % Introduction du chapitre : Expliquer que ce chapitre présente les résultats
% % de l'application du protocole défini au chapitre 4. C'est le "quoi" de l'analyse.

% \section{Application de la Méthode d'Agrégation sur les Portefeuilles de Test}
% % Appliquer la méthode choisie en 4.3 sur les portefeuilles définis en 4.2.

% \subsection{Analyse des Portefeuilles Agrégés}
% % Analyser brièvement les portefeuilles agrégés.
% % Question clé à répondre ici : Les chocs et modifications sont-ils toujours
% % visibles et bien représentés après l'agrégation ?

% \section{Analyse des Sensibilités sur les Indicateurs Solvabilité 2}
% % C'est le cœur de vos résultats.
% % Présenter les résultats issus du modèle ALM pour chaque scénario de sensibilité.
% % Comparer les indicateurs (BE, SCR) entre le portefeuille de base et les portefeuilles modifiés.

% \section{Interprétation des Résultats et Validation de l'Approche}
% % Prendre du recul sur les résultats chiffrés.

% \subsection{Évaluation de la performance du générateur}
% % Dans quelle mesure les résultats confirment-ils que le générateur produit
% % des portefeuilles réactifs et cohérents ?

% \subsection{Synthèse des impacts observés}
% % Quelles sont les différences clés observées ? Quels sont les enseignements
% % principaux des tests de sensibilité ?

\chapter{Conclusion}
\section{Résumé des résultats}
\subsection{Synthèse des principaux résultats obtenus}
\subsection{Impact des méthodes d'agrégation sur les portefeuilles de passifs}

\section{Perspectives d'amélioration}
\subsection{Axes d'amélioration pour les générateurs de portefeuilles de passifs}
\subsection{Évolutions possibles des méthodes d'agrégation et de modélisation ALM}  
\subsection{Autres domaines d'application des générateurs de portefeuilles de passifs}
\section{Conclusion générale}

\chapter*{Annexes}
\addcontentsline{toc}{chapter}{Annexes}
\section*{Annexe chapitre 1}
\section*{Annexe chapitre 2}
\section*{Annexe chapitre 3}
\section*{Annexe chapitre 4}
\section*{Annexe chapitre 5}

\newpage % Pour commencer la bibliographie sur une nouvelle page

% Ajoute une entrée "Bibliographie" dans la table des matières
\addcontentsline{toc}{chapter}{Bibliographie} 

\begin{thebibliography}{9}

% \bibitem{auteur_livre}
% NOM Prénom. \textit{Titre du livre en italique}. Lieu d'édition : Nom de l'éditeur, Année.

% \bibitem{auteur_article}
% NOM Prénom. « Titre de l'article entre guillemets ». \textit{Titre de la revue en italique}, volume X, numéro Y, saison Année, p. 123-456.

\bibitem{clustering_book}
DORNAIKA Fadi, HAMAD Denis, CONSTANTIN Joseph, TRONG HOANG Vinh. \textit{Advances in Data Clustering}. Lieu d'édition : Springer, 2024.

\bibitem{goffard_guerrault}
GOFFARD Pierre-Olivier, GUERRAULT Xavier. « Is it optimal to group policyholders by age, gender, and seniority for BEL computations based on model points? ». \textit{European Actuariel Journal}, volume 5, 17 Avril 2015, p. 165-180.


\bibitem{memoire_ben_fadhel}
BEN FADHEL Amine. « Accéleration de l'évaluation de la solvabilité prospective d'un assureur épargne ». \textit{Mémoire pour l'Institut des Actuaires }, 2022.

\bibitem{france_assureurs}
FRANCE ASSUREURS. \textit{L'assurance vie en 2023}. (20 septembre 2024). Consulté le 1er Septembre 2025, sur \url{https://www.franceassureurs.fr/nos-chiffres-cles/assurance-vie/etude-statistique-assurance-vie-2023/}

% \bibitem{rapport_technique}
% NOM Prénom et NOM Prénom. \textit{Titre du rapport}. Type de rapport (e.g. Rapport de recherche), Organisme/Université, Année.

\end{thebibliography}

\end{document}
